\section*{Глава 1. Анализ предметной области}
\addcontentsline{toc}{section}{Глава 1. Анализ предметной области}
\stepcounter{section}

\subsection{Обзор существующих интернет-магазинов электроники и систем управления заказами}

Рынок цифровой электроники относится к числу наиболее динамичных сегментов электронной коммерции. Покупатели ожидают от интернет-магазина не только удобного каталога товаров, но и подробных технических характеристик, сравнения моделей, актуальной информации о наличии, понятного оформления заказа, выбора способа доставки и прозрачного отслеживания статуса покупки. Для администрации магазина важны другие задачи: управление ассортиментом, контроль остатков, обработка заказов, обновление цен, работа с категориями, пользователями и историей продаж.

В качестве аналогов были рассмотрены крупные интернет-магазины и платформы, работающие с цифровой электроникой: DNS, Ситилинк и Ozon. Эти решения позволяют выявить типовые функциональные возможности, которые целесообразно учитывать при проектировании собственного веб-приложения.

\textbf{DNS} --- российская торговая сеть и интернет-магазин цифровой и бытовой техники. Каталог DNS содержит смартфоны, ноутбуки, комплектующие, периферию, бытовую электронику и аксессуары. Сильной стороной платформы является развитая структура категорий, фильтры по техническим характеристикам и карточки товаров с описанием, ценой, наличием в магазинах, отзывами и вариантами доставки. Для покупателя доступны корзина, оформление заказа, личный кабинет и история покупок. Однако внутренняя система управления заказами ориентирована на процессы самой сети и не может быть использована как самостоятельный инструмент для небольшого магазина электроники.

\textbf{Ситилинк} представляет собой крупный интернет-магазин электроники, компьютерной техники и товаров для дома. Платформа поддерживает поиск по каталогу, подбор товаров по характеристикам, сравнение моделей, оформление доставки или самовывоза, а также работу с корпоративными клиентами. Система хорошо демонстрирует важность точного описания характеристик товара: для цифровой электроники покупатель часто принимает решение на основании параметров процессора, объема памяти, диагонали экрана, емкости аккумулятора и совместимости аксессуаров. Ограничением с точки зрения проектируемого решения является закрытость бизнес-логики: магазин не предоставляет сторонним продавцам гибкий модуль управления собственным ассортиментом и заказами.

\textbf{Ozon} является универсальным маркетплейсом, где цифровая электроника представлена как одна из крупных товарных категорий. Платформа предоставляет покупателю развитые инструменты поиска, фильтрации, сортировки, просмотра отзывов и отслеживания доставки. Для продавцов существует личный кабинет с управлением товарами, ценами, остатками и заказами. При этом работа через маркетплейс предполагает зависимость от правил площадки, комиссий, ограничений интерфейса и общей логики обработки заказов. Для небольшого магазина, которому требуется самостоятельный сайт с собственной базой клиентов и управлением продажами, такой подход не всегда является оптимальным.

Сравнительная характеристика рассмотренных решений приведена в таблице~\ref{tab:electronics_competitors}.

\begin{table}[H]
    \centering
    \caption{Сравнительный анализ интернет-магазинов цифровой электроники}
    \label{tab:electronics_competitors}
    \small
    \begin{tabular}{|p{2.4cm}|p{3.1cm}|p{3.4cm}|p{3.1cm}|p{3.1cm}|}
        \hline
        \textbf{Название} & \textbf{Ассортимент} & \textbf{Управление заказами} & \textbf{Личный кабинет} & \textbf{Гибкость для малого бизнеса} \\ \hline
        DNS & Смартфоны, ноутбуки, комплектующие, периферия, бытовая техника & Оформление покупки, выбор доставки или самовывоза, отслеживание статуса & История заказов, избранное, данные покупателя & Система закрыта и предназначена для собственной торговой сети \\ \hline
        Ситилинк & Электроника, компьютерная техника, комплектующие, товары для офиса и дома & Обработка заказов, самовывоз, доставка, работа с юридическими лицами & Заказы, бонусы, профиль, корпоративные функции & Готового решения для стороннего магазина не предоставляет \\ \hline
        Ozon & Универсальный маркетплейс, включая цифровую электронику & Централизованная обработка заказов через инфраструктуру маркетплейса & История покупок, возвраты, отслеживание доставки & Возможна продажа через площадку, но с зависимостью от правил и комиссий маркетплейса \\ \hline
    \end{tabular}
\end{table}

Проведенный обзор показывает, что существующие крупные платформы обладают развитым пользовательским функционалом, однако их системы управления ассортиментом и заказами либо закрыты, либо жестко связаны с инфраструктурой маркетплейса. Поэтому разработка отдельного веб-приложения для магазина цифровой электроники является актуальной: такое приложение должно объединять публичный каталог товаров, корзину, оформление заказов, личный кабинет покупателя и административный интерфейс для управления товарами, категориями и заказами.

\subsection{Анализ программных средств разработки веб-приложений}

Для реализации веб-приложения необходимо выбрать технологии как для серверной (backend), так и для клиентской (frontend) части.

Для серверной части веб-приложения интернет-магазина цифровой электроники были рассмотрены три популярных Python-фреймворка: Flask, Django и FastAPI. Выбор Python обусловлен простотой разработки, развитой экосистемой библиотек, удобной работой с базами данных и широким применением в учебных и коммерческих веб-проектах.

\textbf{Flask} --- микрофреймворк, предоставляющий базовые инструменты для обработки HTTP-запросов, маршрутизации и формирования ответов. Его преимуществом является простота и гибкость: разработчик самостоятельно выбирает ORM, систему авторизации, валидацию данных и структуру проекта. Такой подход удобен для небольших приложений и прототипов. Недостатком Flask является необходимость вручную подключать и согласовывать многие компоненты, которые в интернет-магазине требуются практически всегда: регистрацию пользователей, разграничение прав, административный интерфейс, работу с формами и защиту операций.

\textbf{Django} --- полнофункциональный веб-фреймворк, ориентированный на быструю разработку надежных серверных приложений. Он включает ORM для работы с базой данных, систему миграций, встроенную административную панель, механизмы аутентификации, работу с формами и шаблонами. Для проекта интернет-магазина эти возможности особенно важны, поскольку требуется хранить товары, категории, пользователей, корзины и заказы, а также предоставить администратору удобный интерфейс управления данными. Django уменьшает объем вспомогательного кода и позволяет сосредоточиться на предметной логике магазина.

\textbf{FastAPI} --- современный фреймворк для создания API, основанный на типизации Python и автоматической генерации документации OpenAPI. Он хорошо подходит для высокопроизводительных REST API и асинхронной обработки запросов. FastAPI удобен, если проект строится как отдельный backend для SPA-приложения или мобильного клиента. При этом для классического интернет-магазина с административной панелью и большим количеством стандартных CRUD-операций часть инфраструктуры потребуется подключать и настраивать отдельно.

Сравнение фреймворков по ключевым критериям приведено в таблице~\ref{tab:framework_compare}.

\begin{table}[H]
    \centering
    \caption{Сравнение Flask, Django и FastAPI}
    \label{tab:framework_compare}
    \small
    \begin{tabular}{|p{3cm}|p{3.6cm}|p{3.6cm}|p{3.6cm}|}
        \hline
        \textbf{Критерий} & \textbf{Flask} & \textbf{Django (выбран)} & \textbf{FastAPI} \\ \hline
        Архитектура & Минимальное ядро, высокая гибкость & Полнофункциональный фреймворк с готовой структурой & API-ориентированный подход с поддержкой асинхронности \\ \hline
        Работа с базой данных & Требуется подключение SQLAlchemy или другого ORM & Встроенная ORM и миграции & ORM выбирается отдельно, часто используется SQLAlchemy \\ \hline
        Администрирование & Готовая панель отсутствует & Встроенная административная панель & Административный интерфейс требуется реализовывать отдельно \\ \hline
        Аутентификация и права доступа & Подключаются расширениями & Встроенные пользователи, группы и права & Реализуются через дополнительные библиотеки и собственную логику \\ \hline
        Подходящая область применения & Небольшие сервисы и прототипы & Интернет-магазины, CRM, панели управления, CRUD-системы & Высоконагруженные API и сервисы для SPA/мобильных клиентов \\ \hline
    \end{tabular}
\end{table}

На основании сравнения для реализации проекта выбран Django. Его встроенные средства позволяют быстрее разработать основные модули интернет-магазина: каталог товаров, категории, пользователей, корзину, оформление заказа и административное управление. Кроме того, Django хорошо подходит для проектов, где важны целостность данных, разграничение доступа и поддерживаемая структура приложения.

В качестве базы данных целесообразно использовать реляционную СУБД SQLite на этапе разработки и PostgreSQL при последующем развертывании. Реляционная модель удобна для описания связей между товарами, категориями, пользователями, заказами и позициями заказа. Для обмена данными между клиентской и серверной частью может применяться формат JSON, особенно при реализации REST API для каталога, корзины и заказов.

Для реализации клиентской части (frontend) были рассмотрены следующие технологии: классическая связка HTML/CSS/JavaScript с использованием шаблонизаторов, а также современные JavaScript-фреймворки (React, Vue.js).

\textbf{HTML/CSS/JavaScript (с шаблонами Django)} --- традиционный подход, при котором сервер генерирует готовые HTML-страницы. Этот вариант отличается простотой интеграции с Django, быстрой первоначальной загрузкой и хорошей SEO-оптимизацией. Дополнительно можно использовать CSS-фреймворки (например, Bootstrap или Tailwind CSS) для ускорения вёрстки адаптивного интерфейса. 

\textbf{React} --- популярная библиотека для создания пользовательских интерфейсов на основе компонентов. React отлично подходит для разработки сложных интерактивных Single Page Application (SPA), обеспечивая высокую скорость отклика без перезагрузки страниц. Однако его использование требует настройки отдельного фронтенд-проекта, сборщика модулей (Webpack/Vite) и организации обмена данными с сервером через REST API.

\textbf{Vue.js} --- прогрессивный JavaScript-фреймворк, который можно использовать как для создания полноценных SPA, так и для постепенного внедрения интерактивности на отдельные страницы. Он обладает более низким порогом вхождения по сравнению с React и предлагает удобную систему реактивности, но также требует дополнительных трудозатрат на настройку API-взаимодействия.

Сравнение технологий клиентской части приведено в таблице~\ref{tab:frontend_compare}.

\begin{table}[H]
    \centering
    \caption{Сравнение технологий клиентской части}
    \label{tab:frontend_compare}
    \small
    \begin{tabular}{|p{3cm}|p{3.6cm}|p{3.6cm}|p{3.6cm}|}
        \hline
        \textbf{Критерий} & \textbf{Шаблоны Django + JS} & \textbf{React} & \textbf{Vue.js} \\ \hline
        Архитектура & Многостраничное приложение (MPA) & Одностраничное приложение (SPA) & SPA или внедрение в MPA \\ \hline
        Сложность интеграции с Django & Минимальная (встроено во фреймворк) & Высокая (требуется REST API и отдельная сборка) & Средняя (возможна постепенная интеграция) \\ \hline
        Интерактивность & Базовая (обновления через перезагрузку страниц или AJAX) & Высокая (динамическое обновление компонентов) & Высокая \\ \hline
        Подходящая область применения & Информационные сайты, простые магазины, админ-панели & Сложные интерактивные веб-приложения и сервисы & Гибкие проекты, где важен баланс между простотой и реактивностью \\ \hline
    \end{tabular}
\end{table}

На основании сравнения, для реализации пользовательского интерфейса интернет-магазина выбран подход с использованием \textbf{HTML/CSS/JavaScript и шаблонизатора Django}, дополненный CSS-фреймворком Bootstrap для адаптивной вёрстки. Этот выбор позволит сократить время на разработку базового функционала магазина (каталог, корзина, оформление заказа) и избежать избыточной сложности, связанной с поддержкой двух отдельных кодовых баз (frontend и backend), что оптимально подходит для масштаба курсового проекта.

\subsection{Формулировка цели, задач и требований к проекту}

На основе анализа предметной области и рассмотренных программных средств была сформулирована цель курсовой работы.

\textbf{Цель работы} --- разработать веб-приложение для управления ассортиментом и заказами магазина цифровой электроники, обеспечивающее покупателям удобную работу с каталогом и заказами, а администраторам --- инструменты ведения товаров, категорий и обработки заявок.

Для достижения поставленной цели необходимо решить следующие \textbf{задачи}:
\begin{enumerate}
    \item Провести анализ существующих интернет-магазинов электроники и определить основные функции, необходимые для проектируемого приложения.
    \item Выбрать программные средства разработки серверной части веб-приложения с учетом требований к каталогу, заказам, пользователям и административному управлению.
    \item Спроектировать структуру данных для хранения товаров, категорий, пользователей, корзины, заказов и позиций заказа.
    \item Реализовать пользовательский интерфейс для просмотра каталога, поиска, фильтрации, добавления товаров в корзину и оформления заказа.
    \item Реализовать аутентификацию пользователей и разграничение ролей покупателей, менеджеров и администраторов.
    \item Создать CRUD-интерфейс для управления товарами, категориями и заказами.
    \item Провести тестирование основных сценариев работы приложения и проверить корректность обработки пользовательских действий.
\end{enumerate}

К проектируемому веб-приложению предъявляются следующие \textbf{функциональные требования}:
\begin{itemize}
    \item отображение каталога цифровой электроники с разделением товаров по категориям;
    \item хранение карточек товаров с названием, описанием, ценой, изображением, характеристиками и количеством на складе;
    \item поиск, фильтрация и сортировка товаров;
    \item регистрация и авторизация пользователей;
    \item добавление товаров в корзину и изменение количества позиций;
    \item оформление заказа с сохранением состава покупки и контактных данных;
    \item просмотр пользователем истории своих заказов;
    \item управление товарами, категориями и заказами через административный интерфейс;
    \item разграничение прав доступа для покупателей, менеджеров и администраторов.
\end{itemize}

Также необходимо учитывать \textbf{нефункциональные требования}: интерфейс должен быть понятным для покупателей без специальной подготовки, данные заказов должны сохраняться в структурированном виде, действия администратора должны быть защищены авторизацией, а архитектура приложения должна допускать дальнейшее расширение, например добавление промокодов, отзывов, интеграции с платежной системой или службами доставки.
